\chapter{当前问题与研究方案调整}\label{chap:middle-discussion}
\section{从剖面可辨识性到真实目标关联}

本阶段磁感知结果说明，主动形成空间变化可以使原本依赖幅值的探测问题获得额外几何约束。半高交点、空间支撑和形状残差将反演条件显式化，使失败能够对应到扫描不足、信号退化或模型不匹配等具体原因。与此同时，场形失配对照揭示，清晰对称的剖面也可能给出稳定而错误的深度。因而问题的重点已经从“如何在理想曲线上计算宽度”推进到“何时允许把宽度解释为某个物理量”。

现有拒绝机制对已声明的模型族建立了可检查的门槛，但未知多芯结构、复合背景和多磁源叠加可能超出该模型集合。增加横切幅度通常有助于观测到剖面两侧，却不能保证唯一辨识所有场形；重复同一种几何的横切也不必然增加独立信息。后续标定需要同时控制相对几何、离底高度、电流或等效激励以及探头姿态，并把不同高度和重复扫描作为检验模型一致性的手段。该过程应优先确定哪些量能够稳定估计，继而决定是否开放深度发布。

目标身份是另一项独立约束。多个来源可以同时指向同一错误位置，融合后方差变小也不意味着目标关联正确。本阶段使用外部受控参考和水平导航误差界的仿真，说明独立几何约束能够影响错误目标授权；这些输入在真实任务中需要由独立调查、可核验目标记录和导航误差评估提供。后续工作必须明确参考信息的来源、适用路线与失效条件，避免把仿真夹具提供的真值便利当成艇载感知已经取得的能力。

埋深解释还受到高度基准的影响。传感器到电缆的垂向距离与海床以下埋深之间，需要通过传感器离底高度连接。若高度、杆臂、姿态或海床起伏的误差未纳入，单一 FWHM 数字的高重复性不能代表最终埋深准确度。因此，在完整标定形成前，报告和运行输出保留深度为空的状态，其含义是该物理量尚无充分依据；可用的横向线索则按自己的条件继续使用。

\section{估计、运动与部署之间的证据缺口}

融合估计实验显示，异常观测处理能够改变误差表现，但滤波结果需要同时接受精度与统计一致性的检验。当前固定协方差、鲁棒权重和时间误差膨胀的比较，并没有形成所有场景下统一占优的结论。后续应在真实传感器时间戳、误差相关性和杆臂条件下核查创新统计，重点区分观测本身的偏置、模型噪声设置和目标关联错误。这些因素若被合并为一个“低置信度”标签，控制层就难以选择合适的响应。

控制侧的主要收益也具有明确职责。UA-MPC 的保守调制能够减小输入变化，但当前代表对照中横向误差没有同步改善，极端场景仍存在约束违规。该结果支持将不确定性用于调节动作强度，同时要求航迹、速度、近底距离和退出策略另有约束。下一步不宜仅通过继续增大惩罚权重追求更平滑的曲线，而应先在实际执行机构上辨识响应、死区与低速转向能力，确定控制器能够利用的运动范围。

运动仿真到软件运行的迁移，还增加了时间与资源因素。Jetson 独立求解器测试为计算负荷提供参考，完整任务却包含传感接收、几何处理、状态发布、记录和控制等并发工作。本机以历史位姿开环重放、重新合成磁测和声学输入并由独立 DDS 进程记录的短试所暴露的连续输出问题，需要在目标平台的同次输入、决策、命令与接收记录中重新定位。单个模块的较小平均耗时、某段仿真的路径到端，均无法覆盖整栈的最长时延和故障传播过程。

实物子链已经提供磁采集、相对几何拟合、计算与真实通信的基础，但整艇装配后还会出现新的耦合。电源和执行设备可能改变磁背景，机械安装可能改变外参，封舱温升可能改变漂移与计算状态，水中负载又会改变执行响应。已有子链结果可以用于设计检查项和选取参数起点；后续整艇试验应重新确认这些条件，而不是把分时、分设备的通过项直接相加。

\section{实施顺序调整与阶段判断}

研究实施已由先行完成算法和接口验证，转向以基础平台恢复牵引实物集成。前一阶段通过独立采集链、仿真与通信子链降低了各模块的不确定性，当前主要缺口集中在它们是否能在同一航行器、同一运行时段和同一任务条件下协同工作。继续累积类似仿真场景对回答这一问题的边际作用有限，后续资源因此优先投入能源、基础反馈和执行链确认，再推进磁载荷装艇与任务验证。

这一调整保留了原课题的研究目标，同时改变了验证顺序。首阶段是否让 Jetson 进入实际执行链，取决于届时确认的基础运行架构：若先利用原控制链恢复受控运动，则新增计算平台可以先承担记录和影子计算；若需要由 Jetson 生成设定值，则其接管、撤回和失联响应必须纳入该阶段准入。不能机械地把新增计算平台全部后置，也不能在没有必要时让复杂任务软件阻塞基础恢复。

内部里程碑 M2.5 指基础平台恢复与水中受控运行验证，涵盖能源、通信、基础反馈、机械与水密、可控执行，以及受控水域中的运动和退出记录。目前缺少完整证据包，尚未完成；后续整艇海缆任务仍须单独验证。
